\chapter{The Missing Obligation: A Refinement Reading}
\label{ch:refinement}
\draftnote{Target length: about 5 pages. Personal contribution; get explicit mentor
approval of the framing before drafting.}

\section{The Problem Restated}
\label{sec:ref-setting}
\draftnote{An abstract policy model $A$, a concrete enforcement model $C$, and a retrieve
relation $\retr$ between their state spaces: the standard setting
\cite{woodcock1996usingz}, applied to access control.}

\section{The Obligations That Would Have to Be Discharged}
\label{sec:ref-obligations}
\draftnote{Initialisation, applicability and correctness (simulation) conditions, written
out concretely for $\tau$ \cite{sissodiya2025k8s} and for the HLP$\to$MLP$\to$LLC
translator \cite{basile2026scm}, so it is clear exactly what is currently being assumed.}

\section{$\Delta(u)$ as a Near-Miss}
\label{sec:ref-delta}
\draftnote{Why the privilege delta in Equation~\eqref{eq:k8s-delta} has the shape of a
retrieve relation, what it lacks to become one, and what would follow if it had it.}

\section{Why Modularity Matters Here}
\label{sec:ref-modularity}
\draftnote{Real policy models are large, and a monolithic obligation does not scale. The
case for refining a small kernel of state-changing operations while leaving derived
operations untouched \cite{faitelson2017decomposition,shepetovsky2026thesis}, advanced as a
proposal, not as a result.}

\section{Honest Limits}
\label{sec:ref-limits}
\draftnote{A refinement proof does not catch a wrong requirement; it catches an unfaithful
implementation of a right one. What this line of work would and would not buy.}
